

\section{L'enseignement comme rituel social : Le grec comme instrument de classement}
space{0.3cm}
oindent

L'analyse sociologique de l'éducation, telle qu'elle a été refondée par Pierre Bourdieu et Jean-
Claude Passeron dans la seconde moitié du XXe siècle, ne se contente pas de corréler des résultats
scolaires à des variables socio-économiques. Elle interroge la nature même des savoirs enseignés,
leur hiérarchie, et la manière dont leur mode de transmission opère une sélection silencieuse mais
implacable. Au cœur de cette critique réside une formule lapidaire mais d'une densité théorique
exceptionnelle, identifiée dans la section I.2.\cite{I-2-3-ref3} de l'architecture conceptuelle
bourdieusienne : \textit{\og L'enseignement comme rituel social : le grec ne sert plus à lire mais à
classer, la difficulté grammaticale maintenue pour trier les 'héritiers'\fg{}}.

Cette assertion, loin d'être une simple critique pédagogique, constitue la clé de voûte d'une
théorie générale de la violence symbolique et de la reproduction sociale. Elle pose une question
fondamentale : comment une discipline académique, les humanités classiques, a-t-elle pu survivre à
l'effondrement de son utilité pratique pour devenir l'instrument suprême de la distinction sociale?
Pour y répondre, il ne suffit pas d'observer les salles de classe des années 1960. Il faut remonter
le fil d'une histoire intellectuelle et institutionnelle complexe, croisant la sociologie de la
reproduction 1, l'histoire des disciplines scolaires théorisée par André Chervel 3, et
l'anthropologie historique de la culture lettrée décrite par Françoise Waquet.\cite{I-2-3-ref5}

Nous
examinerons d'abord comment l'école, sous couvert de neutralité méritocratique, opère une
transmutation de l'héritage social en don naturel. Nous analyserons ensuite la mutation historique
de l'enseignement du grec, passant d'un outil d'accès aux textes à un code de reconnaissance
sociale, un \og signe\fg{} vide de sens mais saturé de valeur. Nous plongerons dans les débats
féroces du début du XXe siècle, notamment autour de la réforme de 1902 et de la méthode directe,
pour démontrer que la difficulté grammaticale fut politiquement et pédagogiquement
\textit{maintenue} pour préserver l'élitisme. Enfin, nous étudierons le mécanisme du concours comme
rituel d'institution, avant de voir comment cette logique de classement a survécu à la mort des
humanités en se transférant vers les mathématiques modernes.



\subsection{La sociologie du privilège : de l'héritage culturel à l'idéologie du don}

Pour comprendre pourquoi \og le grec sert à classer\fg{}, il est impératif de déconstruire le mythe
fondateur de l'école républicaine : l'idéologie du don. C'est dans \textit{Les Héritiers} (1964) que
Bourdieu et Passeron établissent la rupture épistémologique majeure avec la sociologie spontanée des
enseignants.



\subsection{L'illusion de la neutralité scolaire et la réalité statistique}

L'école se vit et se proclame comme une instance neutre, distribuant les places en fonction du seul
mérite individuel. Or, l'analyse statistique révèle une distorsion massive des probabilités d'accès
aux échelons supérieurs de la culture. Les données mobilisées par Bourdieu et Passeron, et
réactualisées par les enquêtes ultérieures, sont sans appel.

\begin{table}[h!]\small\centering\begin{tabular}{| l | l | l | l | l |}\hline\textbf{Catégorie Socio-Professionnelle du Père} & \textbf{Probabilité d'accès (1961-1962)} & \textbf{Ratio de chance (Base 1 \= Ouvrier)} & \textbf{Probabilité d'accès (2005-2006)} & \textbf{Ratio de chance (Base 1 \= Ouvrier)} \\ \hline\textbf{Ouvriers / Salariés agricoles} & \~1-2 \% & \textbf{1} & \~20 \% & \textbf{1} \\ \hline\textbf{Cadres supérieurs / Professions libérales} & \> 40 \% & \textbf{\~42} & \> 75 \% & \textbf{\~4} \\ \hline\end{tabular}\end{table}\textit{Note : Les ratios indiquent qu'en 1961, un fils de cadre avait 42 fois plus de chances
d'entrer à l'université qu'un fils d'ouvrier. Bien que cet écart se soit réduit quantitativement en
2005 (ratio de 4), la ségrégation s'est déplacée vers la nature des filières (Grandes Écoles vs
Université de masse).} 1

Cette inégalité objective ne peut s'expliquer par une inégalité de nature (le \og don\fg{}). Elle
s'explique par la rentabilisation scolaire d'un capital culturel familial. L'école ne transmet pas
la culture ; elle sanctionne ceux qui la possèdent déjà. Elle exige, dès l'entrée, une familiarité
avec le langage, une aisance rhétorique, et une posture corporelle (\og l'hexis\fg{}) que seule la
famille bourgeoise transmet par osmose quotidienne.\cite{I-2-3-ref1}



\subsection{Le grec ancien comme capital culturel objectivé}

Dans cette économie des échanges symboliques, le grec ancien occupe une place singulière. Il ne
s'agit pas d'un savoir utile (comme l'anglais ou la comptabilité), ni d'un savoir civique (comme
l'histoire de France). Il s'agit d'un savoir de \textit{luxe}.

\begin{itemize}\item \textbf{La gratuité comme valeur :} La sociologie de la distinction, développée ultérieurement par Bourdieu, postule que plus une pratique est éloignée de la nécessité économique, plus elle est distinctive.\cite{I-2-3-ref6} Le grec, par son apparente inutilité marchande, devient le marqueur suprême du désintéressement bourgeois. Étudier le grec, c'est prouver que l'on a le temps (\og skholè\fg{}) de ne rien faire d'utile.\item \textbf{L'affinité élective :} L'enseignement classique valorise des qualités — l'élégance, la finesse, le \og sentir\fg{} — qui sont précisément les attributs de la conversation mondaine. L'héritier n'a pas besoin de \og bûcher\fg{} sa grammaire grecque avec l'acharnement du boursier ; il aborde le texte avec une familiarité désinvolte, reconnaissant dans Homère ou Platon les références qui peuplent déjà la bibliothèque paternelle.\cite{I-2-3-ref8}\end{itemize}

\subsection{Le mécanisme de la méconnaissance}

La force du système réside dans sa capacité à dissimuler sa fonction sociale. Si l'école affichait
ouvertement \og nous sélectionnons les riches\fg{}, elle perdrait sa légitimité. Elle doit donc dire
\og nous sélectionnons les meilleurs\fg{}.

L'enseignement du grec sert de masque parfait. La difficulté technique (déclinaisons, verbes
irréguliers) fournit un alibi objectif à l'élimination. L'élève qui échoue en version grecque est
convaincu de sa propre insuffisance intellectuelle (\og je ne suis pas doué pour les langues\fg{}).
Il intériorise son échec social comme un échec personnel, validant ainsi la hiérarchie qui l'exclut.
C'est ce que Bourdieu nomme la violence symbolique : une violence qui s'exerce avec la complicité
tacite de ceux qui la subissent.\cite{I-2-3-ref10}



\subsection{Archéologie d'une discipline : de la lecture au classement}

L'assertion \og le grec ne sert plus à lire\fg{} n'est pas une simple métaphore. Elle correspond à
une réalité historique documentée par l'histoire des disciplines scolaires. André Chervel et
Françoise Waquet ont montré comment, au fil des siècles, la finalité de l'enseignement des langues
anciennes a basculé.



\subsection{L'empire du signe : le latin et le grec comme valeurs refuges}

Françoise Waquet, dans son ouvrage magistral \textit{Le latin ou l'empire d'un signe}, démontre que
dès le XVIIIe siècle, et massivement au XIXe, la compétence réelle des élèves en latin et en grec
s'effondre.

\begin{itemize}\item \textbf{Le constat d'échec :} Les rapports d'inspection et les témoignages d'époque concordent : après sept ou huit années d'études acharnées, la vaste majorité des bacheliers est incapable de lire un texte de Cicéron ou de Platon à livre ouvert.\cite{I-2-3-ref5} Si la fonction de l'enseignement était l'apprentissage de la langue (fonction manifeste), le système aurait dû être déclaré en faillite.\item \textbf{La réussite du signe :} Si le système perdure, c'est que sa fonction latente est remplie. Le latin et le grec ne valent plus comme langues de communication (ce qu'ils étaient encore à la Renaissance ou dans l'Église), mais comme \textit{signes}. Le simple fait d'avoir \og fait ses humanités\fg{} imprime sur l'élève une marque indélébile d'appartenance à l'élite. Le contenu s'efface devant le contenant. Le grec devient un \og monument\fg{} que l'on visite rituellement, non pour l'habiter, mais pour en revenir sanctifié.\cite{I-2-3-ref5}\end{itemize}

\subsection{La mutation des exercices : l'hégémonie de la version}

Cette transformation fonctionnelle se lit dans l'évolution matérielle des exercices scolaires
étudiée par André Chervel.\cite{I-2-3-ref3}

Jusqu'au XVIIIe siècle, l'exercice roi était le thème (écrire en latin/grec) ou le discours
(rhétorique). On visait une production active. Au XIXe siècle, une révolution silencieuse s'opère :
la version (traduction vers le français) devient l'exercice dominant, puis exclusif, des examens
(Baccalauréat, Agrégation).

\begin{table}[h!]\small\centering\begin{tabular}{| p{2cm} | p{2.2cm} | p{2.2cm} | p{2.2cm} |}\hline\textbf{Période} & \textbf{Exercice Dominant} & \textbf{Compétence Visée} & \textbf{Fonction Sociale} \\ \hline\textbf{Ancien Régime} & Thème, Vers latins, Discours & Production, Éloquence, Imitation active & Formation du Clergé et de la Robe (usage professionnel du latin) \\ \hline\textbf{XIXe siècle} & Version écrite & Compréhension passive, Maîtrise du Français & Formation de la Bourgeoisie (usage mondain de la culture) \\ \hline\textbf{XXe siècle} & Version \og sèche\fg{} (extrait court) & Analyse grammaticale, \og Gymnastique mentale\fg{} & Sélection républicaine (classement par la difficulté formelle) \\ \hline\end{tabular}\end{table}Pourquoi ce basculement? Parce que la version permet de classer plus finement.

\begin{itemize}\item \textbf{La version comme exercice de français :} Paradoxalement, la version grecque teste moins la connaissance du grec que la maîtrise du français littéraire.\cite{I-2-3-ref14} Pour traduire \og bien\fg{}, il ne suffit pas de comprendre le sens (ce que le \og boursier\fg{} laborieux parvient à faire avec son dictionnaire). Il faut trouver le \og mot juste\fg{}, la tournure élégante, le subjonctif imparfait qui \og sonne\fg{} bien.\item \textbf{Le piège du style :} C'est ici que l'héritier triomphe. Sa traduction est \og élégante\fg{}, celle du parvenu est \og lourde\fg{}, \og scolaire\fg{}, ou \og servile\fg{}. Le vocabulaire de la notation (\og plat\fg{}, \og jargon\fg{}, \og manque de noblesse\fg{}) est un vocabulaire esthétique qui masque un jugement de classe.\cite{I-2-3-ref15} Le grec ne sert plus à accéder à la pensée de Platon (souvent déjà connue par ailleurs), mais à fournir un prétexte pour évaluer la distinction linguistique en français.\end{itemize}

\subsection{L'abandon de la lecture cursive}

La conséquence pédagogique de cette logique de classement est l'abandon de la lecture des œuvres
complètes. On ne lit plus L'Iliade. On traduit vingt lignes du chant VI, choisies pour leurs
difficultés syntaxiques.

Le texte est réduit à l'état de \og prétexte grammatical\fg{}.\cite{I-2-3-ref16} Il est découpé,
extrait de son contexte, transformé en champ de mines pour l'évaluation. Cette pratique interdit
l'accès au sens et au plaisir du texte, réservant ce plaisir à ceux qui peuvent le lire par ailleurs
(la culture libre de l'héritier). Pour l'élève scolaire, le texte grec est un cadavre que l'on
dissèque, pas une voix que l'on écoute.



\subsection{La politique de la difficulté : les guerres scolaires (1902-1923)}

L'idée que la difficulté grammaticale est \textit{maintenue} volontairement n'est pas une vue de
l'esprit conspirationniste. Elle est documentée par les débats politiques et pédagogiques féroces
qui ont secoué la Troisième République, cristallisant l'opposition entre les \og Modernes\fg{} et
les \og Classiques\fg{}.



\subsection{Le traumatisme de 1902 et la "crise du français"}

La réforme de 1902 constitue un tournant majeur. Elle instaure l'égalité de sanction entre le
baccalauréat classique (A) et le baccalauréat moderne (sans latin/grec, section
D).\cite{I-2-3-ref17} Pour la première fois, on reconnaît qu'une culture fondée sur les sciences et
les langues vivantes peut former l'élite.

La réaction des conservateurs est immédiate et violente. On parle de \og crise du français\fg{}, de
\og baisse du niveau\fg{}, de \og barbarie scientifique\fg{}.\cite{I-2-3-ref6} Les partisans des
humanités, menés par des figures comme Léon Bérard (futur ministre) ou Henri Bergson (en coulisses),
mobilisent un argumentaire qui deviendra classique : les humanités ne servent pas à être utilisées,
elles servent à former l'esprit.



\subsection{La controverse de la méthode directe : le refus de la facilité}

Au même moment, un mouvement pédagogique européen tente de rénover l'enseignement des langues. C'est
l'essor de la \og Méthode Directe\fg{} pour les langues vivantes (bain linguistique, oralité,
interdiction de la traduction).\cite{I-2-3-ref19}

Des pédagogues audacieux proposent d'appliquer cette méthode au latin et au grec (méthode de Rouse
ou d'Officiers). L'idée est simple : pour savoir lire le grec, il faut le parler, le vivre, le
rendre vivant et facile d'accès.

Le rejet de cette méthode par l'institution universitaire française est symptomatique.

\begin{itemize}\item \textbf{L'argument de l'effort :} Les détracteurs de la méthode directe soutiennent que si le grec devient facile, il perd sa valeur éducative. On fustige une méthode \og ludique\fg{}, \og superficielle\fg{}, un \og gaspillage d'énergie\fg{} noble au profit d'un utilitarisme vil.\cite{I-2-3-ref21}\item \textbf{La difficulté comme vertu :} On théorise alors la \og gymnastique mentale\fg{}.\cite{I-2-3-ref22} De même que le corps a besoin d'efforts physiques gratuits pour se muscler, l'esprit a besoin de la résistance de la grammaire grecque pour se former. L'aridité, l'ennui, la complexité des règles (verbes en \-mi, optatifs obliques) ne sont pas des défauts pédagogiques à corriger, mais des vertus ascétiques à préserver. Elles garantissent que seuls les esprits capables d'abstraction et de volonté (c'est-à-dire, sociologiquement, les enfants libérés des contraintes matérielles) pourront réussir.\end{itemize}

\subsection{La restauration de 1923 : le "décret bérard"}

L'apogée de cette politique de la barrière est la réforme Bérard de 1923. Léon Bérard, ministre de
l'Instruction publique, supprime purement et simplement la section moderne au lycée. Il rend le
latin et le grec obligatoires pour tous les élèves désirant faire des études secondaires
complètes.\cite{I-2-3-ref17}

\begin{itemize}\item \textbf{L'objectif de classe :} L'objectif, à peine voilé, est de freiner l'afflux des enfants des classes moyennes et populaires (issus des Écoles Primaires Supérieures) vers le lycée. Ces enfants, souvent excellents en sciences et en français, n'ont pas le bagage latin. En remettant la barrière du latin à l'entrée, on \og trie\fg{} socialement la population scolaire sous couvert d'unité culturelle.\item \textbf{L'échec politique, la réussite symbolique :} Si la réforme est annulée en 1924 par le Cartel des Gauches face au tollé, elle a réussi à ancrer durablement l'idée que la \og vraie\fg{} culture est la culture classique. La filière A (Lettres Classiques) restera jusqu'aux années 1960 la voie royale de l'excellence, celle des \og héritiers\fg{} par excellence.\cite{I-2-3-ref18}\end{itemize}

\subsection{Le rituel de la "botte" : mécanique de la sélection par la grammaire}

Comment la difficulté grammaticale opère-t-elle concrètement le tri dans la salle d'examen? Il faut
entrer dans la \og boîte noire\fg{} de la notation et de la docimologie (science des examens) telle
que Bourdieu l'analyse dans \textit{La Noblesse d'État}.



\subsection{La hiérarchie des fautes et les "points négatifs"}

L'évaluation de la version grecque repose traditionnellement sur un barème de \og points
négatifs\fg{}. On part de 20 et on retranche des points à chaque erreur.\cite{I-2-3-ref24} Mais
toutes les erreurs ne se valent pas.

\begin{itemize}\item \textbf{Le Contresens vs le Faux-sens :} Le système distingue des fautes techniques (le contresens, le barbarisme) qui sont lourdement sanctionnées, et des fautes de goût (le faux-sens léger, la maladresse).\item \textbf{L'insécurité linguistique du petit-bourgeois :} L'élève d'origine modeste, conscient de sa fragilité, développe une obsession de la règle. Il veut être \og correct\fg{}. Cette anxiété le conduit souvent à l'hyper-correction ou à une traduction littérale, servile, \og mot à mot\fg{}, qui est la hantise du correcteur.\cite{I-2-3-ref25} Il connaît ses déclinaisons par cœur, mais il \og ne sent pas\fg{} la langue.\item \textbf{L'aisance de l'héritier :} L'héritier, lui, peut se permettre de prendre des libertés avec la lettre pour sauver l'esprit. Il devine le sens grâce à sa culture générale (il reconnaît l'épisode mythologique). Sa traduction peut comporter des inexactitudes grammaticales mineures, mais elle sera rachetée par le \og brio\fg{} du style français. Le correcteur dira : \og C'est inexact, mais c'est bien dit\fg{}, valorisant l'habitus bourgeois sur la compétence scolaire stricte.\end{itemize}

\subsection{Les exceptions comme pièges à "bûcheurs"}

La grammaire grecque est célèbre pour ses exceptions. Les verbes irréguliers, les dialectes (ionien,
dorien, éolien), les règles d'accentuation complexes (loi de la pénultième) forment un maquis
impénétrable pour le néophyte.

\begin{itemize}\item \textbf{La fonction de l'exception :} Dans une optique de lecture, ces exceptions sont anecdotiques (le contexte suffit souvent à comprendre). Dans une optique de \textit{classement}, elles sont centrales. C'est sur l'exception que l'on piège le \og bûcheur\fg{}. L'élève scolaire a appris la règle ; l'examen porte sur l'exception. Seul celui qui possède une culture surabondante, ou qui a bénéficié d'une préparation spécifique (les précepteurs, les parents professeurs), peut déjouer ces pièges.\cite{I-2-3-ref25}\item \textbf{L'élimination statistique :} En maintenant ces difficultés artificielles (on pourrait simplifier l'enseignement pour la lecture), l'institution s'assure d'une courbe de Gauss des notes très étalée. Cela permet de produire un classement incontestable pour les concours (ENS, Agrégation). La difficulté grammaticale est l'outil technique qui fabrique la rareté des élus.\end{itemize}

\subsection{Le concours comme rite d'institution}

Bourdieu compare le concours (notamment l'Agrégation ou l'entrée à l'ENS) à un rite de passage
anthropologique.\cite{I-2-3-ref28}

\begin{itemize}\item \textbf{La séparation :} Le rite institue une différence de nature. Avant le concours, on est un étudiant ; après, on est un \og agrégé\fg{}, un être à part.\item \textbf{L'épreuve qualifiante :} Pour que cette magie sociale opère, il faut que l'épreuve soit terrible. Il faut que le candidat ait souffert. La difficulté du grec, les nuits blanches, l'ascèse de la préparation (la \og khâgne\fg{}) fonctionnent comme les épreuves d'endurance des rites initiatiques. Elles prouvent que le candidat est digne d'entrer dans la \og Noblesse d'État\fg{}. Si le grec servait à lire, on pourrait l'apprendre joyeusement. Parce qu'il sert à consacrer, il doit être appris dans la douleur.\cite{I-2-3-ref30}\end{itemize}

\subsection{La métamorphose du rituel : du grec aux mathématiques}

L'analyse de Bourdieu ne se fige pas dans la nostalgie. Dans \textit{La Noblesse d'État} (1989) et
dans ses cours ultérieurs, il observe un phénomène fascinant : le transfert de la fonction de
sélection du grec vers les mathématiques modernes.



\subsection{Le déclin des humanités et la relève scientifique}

À partir des années 1960, sous l'effet de la modernisation économique et de la démocratisation
scolaire, le monopole du latin-grec s'effrite. La section A perd de son prestige au profit de la
section C (Scientifique).\cite{I-2-3-ref6}

Est-ce la fin de la sélection sociale? Non, répond Bourdieu. C'est une translation. Les \og
Mathématiques Modernes\fg{} (théorie des ensembles, abstraction algébrique) ont repris le flambeau.



\subsection{Les mathématiques comme nouveau latin}

Les mathématiques partagent avec le grec ancien les caractéristiques essentielles pour servir de
rituel de sélection :

1. \textbf{L'Abstraction :} Elles sont déconnectées de l'expérience immédiate.

2. \textbf{La Sélectivité :} Elles permettent un classement rigoureux (vrai/faux) et hiérarchisé.

3. \textbf{L'Innéisme supposé :} Comme pour le grec (\og la bosse des langues\fg{}), on croit au \og
don\fg{} pour les maths (\og la bosse des maths\fg{}).

4. \textbf{La Gratuité apparente :} Les maths de l'X ou de l'ENS ne sont pas les maths utilitaires
de l'ingénieur ; ce sont des mathématiques pures, théoriques, qui fonctionnent comme un signe
d'excellence intellectuelle générale.\cite{I-2-3-ref1}

La structure du rituel est conservée, seul le contenu a changé. Les \og héritiers\fg{} se sont
massivement reportés vers les filières scientifiques pour maintenir leur position dominante.
Cependant, Bourdieu note une nuance : la maîtrise de la culture littéraire classique reste, pour les
scientifiques dominants, un \og supplément d'âme\fg{}, une marque de distinction ultime qui les
sépare des simples techniciens. Le PDG qui cite Platon en grec (ou qui en a l'air) conserve une aura
supérieure à celui qui ne sait que compter.\cite{I-2-3-ref7}

L'exégèse de la section I.2.\cite{I-2-3-ref3} nous a menés au cœur de la machine scolaire. La thèse
de Bourdieu, selon laquelle \og le grec ne sert plus à lire mais à classer\fg{}, apparaît non comme
une boutade polémique, mais comme le résumé d'un siècle d'histoire éducative.

Cette analyse révèle trois enseignements majeurs pour la sociologie contemporaine :

1. \textbf{La primauté de la fonction latente :} Une discipline peut survivre institutionnellement
bien après sa mort pédagogique ou utilitaire, pourvu qu'elle remplisse une fonction sociale de tri.
L'inefficacité de l'enseignement du grec (l'incapacité à lire) était la condition même de son
efficacité sociale (la capacité à éliminer).

2. \textbf{La construction politique de la difficulté :} La difficulté scolaire n'est pas une donnée
naturelle. Elle est construite, calibrée et maintenue par des choix pédagogiques (préférence pour la
version écrite, rejet de l'oralité, culte de l'exception grammaticale) qui répondent à des
impératifs de régulation sociale.

3. \textbf{L'infinie plasticité du privilège :} Si le grec a disparu comme outil de masse, la
logique de la distinction s'est recomposée ailleurs. Que ce soit à travers les mathématiques pures,
les filières bilingues, ou les \og soft skills\fg{} aujourd'hui, le système éducatif continue de
générer des \og rituels de classement\fg{} qui transforment l'héritage en mérite.

En définitive, étudier l'enseignement du grec comme \og rituel social\fg{}, c'est comprendre que
l'école n'est pas seulement le lieu de la transmission des savoirs, mais le lieu de la production
légitime des hiérarchies sociales. Le grec ancien fut, pendant un siècle, le langage chiffré de la
bourgeoisie française, un mot de passe permettant aux initiés de se reconnaître et d'exclure les
profanes au seuil du temple scolaire.

\begin{itemize}\item \textbf{Bourdieu, P., \& Passeron, J.-C.} (1964). \textit{Les Héritiers : Les étudiants et la culture}. Paris : Les Éditions de Minuit.\cite{I-2-3-ref1}\item \textbf{Bourdieu, P., \& Passeron, J.-C.} (1970). \textit{La Reproduction : Éléments pour une théorie du système d'enseignement}. Paris : Les Éditions de Minuit.\cite{I-2-3-ref11}\item \textbf{Bourdieu, P.} (1979). \textit{La Distinction : Critique sociale du jugement}. Paris : Les Éditions de Minuit.\cite{I-2-3-ref7}\item \textbf{Bourdieu, P.} (1989). \textit{La Noblesse d'État : Grandes écoles et esprit de corps}. Paris : Les Éditions de Minuit.\cite{I-2-3-ref10}\item \textbf{Chervel, A.} (1998). \textit{La culture scolaire. Une approche historique}. Paris : Belin.\cite{I-2-3-ref3}\item \textbf{Waquet, F.} (1998). \textit{Le latin ou l'empire d'un signe, XVIe-XXe siècle}. Paris : Albin Michel.\cite{I-2-3-ref5}\item \textbf{Puren, C.} (1988). \textit{Histoire des méthodologies de l'enseignement des langues}. Paris : Nathan.\cite{I-2-3-ref19}\item \textbf{Jarrety, M.} (2007). \textit{Bergson et la réforme de 1923}. Fabula.\cite{I-2-3-ref17}\end{itemize}